\documentclass[UTF8,a4paper,11pt,fontset=fandol]{ctexart}
\usepackage[margin=23mm]{geometry}
\usepackage{booktabs,tabularx,amsmath}
\usepackage[colorlinks=true,urlcolor=blue]{hyperref}
\setlength{\emergencystretch}{2em}
\title{A20--AEI R9 三视角内部审查}
\author{内部技术核查；非独立外部同行评审}\date{2026年10月12日}
\begin{document}\maketitle
\section{结论}
本轮比只修改文字或增加普通随机几何样本更实质：实际执行三种工程任务接口、生产者变化实验和公平成本对照，但最初17页压缩稿削弱了原有定位研究主线，用户已指出该编辑问题。修正版恢复R8建模、契约、算法、理论和全组定位/航线结果，再整合新证据。它仍不支持预设9分或Strong Accept。若作为投稿候选，本内部判断仍建议\textbf{Major Revision}，核心风险是知识表示原创性、完整服务可用性与跨无关任务一般性。没有邀请外部审稿人，也没有将同一代理的三个检查视角冒称三位独立审查员。

\section{视角一：理论与实现}
事件包含给出点态失败集合包含；请求投影函数不变给出相同事件，进而继承已有锚点上界。逻辑在冻结程序、匹配总体和可信依赖说明下成立。允许条件不能简化为“元数据相同”，不能外推为未知漂移检测、真实设备证明或逐图条件安全。

代码检查保守补入捕获/运动提供者依赖，防止遗漏。故障锁绑定可信来源，错误报文不能靠改来源ID逃离拒绝。接收器是后处理观察器，不能宣传成新的实际相机实时管线。JSON数组恢复到规定元组字段后，完整程序式重新查询与F1一致，序列化没有被用来削弱对手。

\section{视角二：统计与实验}
12,800个独立来源、8,192校准与4,096确认不相交；策略、工况、要求重复不加样本量。新卡统一使用0.05/4,608分配。对全部3,072卡做精确上界核对，独立二项反演检查通过。报告联合事件直接校准；不从两个边际推导同等联合界。

预留族、种子、阈值、预测器和数值服务在确认前冻结。两个校准后修正只增强适用依赖与基线信息公平性，零确认输入时完成并保留旧F1，未改变数值发生器、任务阈值和校准样本。数据确认之后没有挑选有利阈值或策略。

完整关联的一个通过配置仍只有0.006027裕度；完整采集/报告仍拒绝。前缀/置信灵敏度只作诊断，没有额外网格同时保证。五个实际重放与全部原始文件哈希通过，但Linux跨平台全量复现没有做。

\section{视角三：AEI贡献与公平性}
明确工程量、单位、要求变化、证据血缘和依赖闭包，比给定位加一个知识图名词更贴近工程信息学。新增像素采集、因果关联和交付事件，是实际可执行的消费任务扩展。不过它们共享一个几何生产者，不等于跨无关领域一般化。

完整程序式实现匹配契约全部结果，直接续证支持更多配置。少做53.66\%记录评估是锚点语义复用的收益，程序式也有；不是存储形式独有创新。各方法均可缓存重复请求。没有真实人工维护试验，不能声称降低工程师开发成本。原Fixed2优势、原航线和四张真实照片/五张原仿真视图均保留。

\section{仍需新的实质证据}
下一轮优先验证完整任务可用性和真实要求演化，而不堆普通几何样本。改变GSD需作用于焦距、采样尺度或距离，不能用扩大ROI代替；光学与成本修改都需要新生产者和独立续证。另一种无关工程消费任务、更直接的依赖维护/证据迁移评价，以及独立审查，比继续润色同一结果更能降低AEI原创性风险。

所有负面结果与适用边界应继续保留。风险放宽到0.20、删除任务要求、把虚拟渲染时钟叫物理相机延迟，均不能用来虚构完整工程优势。投稿前请实际作者核验声明与责任；本内部报告只提供技术检查与当前局限。
\end{document}
